\begin{textsample}{NEG Dim 1 – 2012\_c5 – Score 6.00 – 2012\_c5/enade\_2012\_admin\_bb7f435184.txt}  \label{ex:f1_neg_327}
CAPÍTULO 5 DISTRIBUIÇÃO DOS CONCEITOS 5.1 PANORAMA NACIONAL DA DISTRIBUIÇÃO DOS CONCEITOS A Tabela 5.1 apresenta a quantidade e distribuição de cursos de Administração participantes do ENADE/2012, por faixa de conceito e Grande Região.
A diferença entre os cursos tabulados neste capítulo e no capítulo 2 são os cursos sem conceito, em princípio, aqueles sem alunos concluintes que participassem da prova ou que tivessem tão somente um aluno concluinte11.
Observando -se os dados da Tabela 5.1, nota -se que, dos 1.554 cursos participantes, 672 ( 43,2 \% ) classificaram -se com conceito 3, o valor modal.
Este foi também o conceito modal em todas as regiões do Brasil ( entre 41,5 \% na região Sudeste e 50,0 \% na região Sul ), exceto na região Centro-Oeste, onde a moda foi o conceito 2 ( 46,2 \% ).
A região Sul apresentou a maior proporção no conceito 4, com 18,5 \% dos cursos com este conceito.
O conceito 2 foi o segundo mais frequente em nível nacional ( 33,7 \%, correspondendo a 524 cursos ) e o conceito 4, o terceiro ( 13,8 \%, correspondendo a 215 cursos ).
Houve, ainda, 96 cursos ( 6,2 \% ) que receberam conceito 5 e outros 33 cursos que receberam conceito 1.
Dos 1.554 cursos de Administração, 14 ficaram sem conceito ( SC ).
A região Norte participou com 88 cursos ou 5,7 \% do total nacional.
Nesta região, o conceito 3 foi atribuído a 42 cursos dos 88 participantes, o que equivale a 47,7 \% do total regional.
Apenas um curso ficou sem conceito ( 1,1 \% ).
Os demais foram avaliados com conceito 1 ( três cursos, correspondendo a 3,4 \% ), conceito 2 ( 35 cursos correspondendo a 39,8 \% ) e conceito 4 ( sete cursos correspondendo a 8,0 \% ). 11 Estes cursos com somente um concluinte foram considerados no capítulo 2. 39,8 \% ) e conceito 4 ( sete cursos correspondendo a 8,0 \% ).
Nenhum curso foi alocado ao conceito 5.
A região Nordeste participou com 262 cursos ou 16,9 \% do total nacional.
Nesta região, 110 cursos, 42,0 \% em termos regionais, obtiveram conceito 3, o conceito modal para a região.
Os conceitos 2 e 4 foram atribuídos a, respectivamente, 94 e 36 cursos ( 35,9 \% e 13,7 \% ).
Já os conceitos 1 e 5 foram atribuídos a, respectivamente, sete e 14 cursos ( 2,7 \% e 5,3 \% ).
Nessa região, apenas um dos cursos ficou sem conceito ( 0,4 \% ).
Como já comentado, dos 697 cursos participantes da região Sudeste, 289 ( 41,5 \% ) obtiveram conceito 3, o conceito modal.
O conceito 1 foi atribuído a 14 cursos ( 2,0 \% ) e o conceito 2, a 233 ( 33,4 \% ).
O conceito 4 foi atribuído a 97 cursos ( 13,9 \% ) e 55 outros ( 7,9 \% ) receberam o conceito 5.
Apenas 9 cursos ficaram sem conceito ( 1,3 \% ).
A região Sul também contou com cursos distribuídos em todas as faixas de conceitos.
A predominância do conceito 3 foi de 50,0 \%, como já comentado anteriormente, correspondentes a 168 dos 336 cursos participantes na região Sul.
O conceito 2 foi atribuído a 83 cursos ( 24,7 \% ) e o conceito 4, a 62 cursos ( 18,5 \% ).
O conceito 1 foi recebido por um único curso ( 0,3 \% ) e o conceito 5, por 19 deles ( 5,7 \% ).
Apenas 3 dos cursos da região Sul ficaram sem conceito ( 0,9 \% ).
Como já comentado, de todos os 171 cursos participantes na região Centro-Oeste, 79 receberam conceito 2, o conceito modal ( 46,2 \% ).
Nesta região nenhum curso ficou sem conceito ( SC ).
Já, oito cursos foram alocados no conceito 1 e outros oito no conceito 5, correspondendo a 4,7 \% dos cursos da região.
Os demais cursos foram avaliados com conceito 3 ( 63 cursos, 36,8 \% ) e conceito 4 ( 13 cursos, 7,6 \% ). 5.2 CONCEITOS POR CATEGORIA ADMINISTRATIVA E POR GRANDE REGIÃO A Tabela 5.2 apresenta a distribuição dos cursos participantes do ENADE/2012 de Administração, por Categoria Administrativa, de acordo com os conceitos por eles alcançados, segundo as Grandes Regiões brasileiras.
Dos 1.554 cursos participantes, 197 ( 12,7 \% ) eram ministrados em instituições públicas e 1.357 ( 87,3 \% ), em privadas.
De acordo com as informações da Tabela 5.2, em termos nacionais, entre instituições públicas e privadas : das 96 IES que receberam o conceito 5, 34 eram públicas e 62 privadas.
Dos 197 cursos participantes de IES públicas, o conceito 3 foi o valor modal, atribuído a 76 cursos ( 38,6 \% ).
Entre os demais cursos de instituições públicas participantes, três obtiveram conceito 1 ( 1,5 \% da categoria ), 22 receberam conceito 2 ( 11,2 \% ) e 61 foram alocados ao conceito 4 ( 31,0 \% ).
Como já comentado, os 34 restantes receberam o conceito 5 ( 17,3 \% ).
Nesta categoria, apenas um dos cursos ficou sem conceito ( SC ).
Na rede privada, o conceito modal também foi 3, com 596 cursos dos 1.357 da categoria.
Entre os demais cursos participantes, 30 receberam conceito 1 e 502, conceito 2.
O conceito 4 foi atribuído a 154 cursos e o conceito 5, como já comentado, a 62 cursos.
Nesta categoria administrativa, 13 dos cursos ficaram sem conceito ( SC ).
Como já comentado, os 34 restantes receberam o conceito 5 ( 17,3 \% ).
Nesta categoria, apenas um dos cursos ficou sem conceito ( SC ).
Na rede privada, o conceito modal também foi 3, com 596 cursos dos 1.357 da categoria.
Entre os demais cursos participantes, 30 receberam conceito 1 e 502, conceito 2.
O conceito 4 foi atribuído a 154 cursos e o conceito 5, como já comentado, a 62 cursos.
Nesta categoria administrativa, 13 dos cursos ficaram sem conceito ( SC ).
Na análise por região, observa -se que, na região Norte, as instituições privadas participaram com 73 cursos ( 83,0 \% do total regional ), dos quais nenhum ficou sem conceito ( SC ).
O conceito modal para as instituições privadas na região foi 3, com 33 cursos, correspondendo a 45,2 \%.
O conceito 2 apresenta frequência muito próxima da do conceito modal, com 32 cursos, correspondendo a 43,8 \%.
Outros três cursos receberam conceito 1 ( 4,1 \% ).
Apenas 5 cursos foram alocados no conceitos 4 ( 6,8 \% ), enquanto nenhum foi alocado no conceito 5.
As instituições públicas participaram com 15 cursos ( 17,0 \% do total regional ), dos quais apenas um ficou sem conceito ( SC ).
O conceito modal para esta Categoria Administrativa, na região Norte, também foi 3, com nove cursos, correspondendo a 60,0 \%.
Os conceitos 2 e 4 participaram com três e dois cursos ( 20,0 \% e 13,3 \% ), respectivamente.
Nenhum curso foi alocado aos demais conceitos nesta categoria na região.
Na região Nordeste, a rede privada concentrou 208 dos 262 cursos participantes, equivalentes a 79,4 \% do total da região, a menor proporção entre todas as regiões.
Nesta Categoria Administrativa na região, predominaram os cursos com conceito 2, 92 cursos correspondendo a 44,2 \% dos cursos oferecidos por IES privadas no Nordeste.
Apenas seis cursos receberam o conceito 1 ( 2,9 \% ), enquanto o conceito 5 participou com sete cursos ( 3,4 \% ).
Os demais foram avaliados com conceitos 3 ( 85 cursos correspondendo a 40,9 \% ) e 4 ( 17 cursos correspondendo a 8,2 \% ).
As instituições públicas dessa região participaram com 54 cursos ( 20,6 \% ), dos quais 25 obtiveram conceitos 3, o valor modal.
Os demais foram avaliados com os conceitos 1 ( um curso ), 2 ( dois cursos ), 4 ( 19 cursos ) e 5 ( sete cursos ).
Nenhum curso ficou sem conceito ( SC ) nesta combinação de região e categoria.
Na região Sudeste, a proporção de cursos da rede privada, 93,0 \%, foi mais elevada do que nas demais regiões brasileiras, correspondendo a 648 dos 697 cursos participantes.
Nesta categoria, na região Sudeste, o conceito modal foi 3 ( 279 cursos ).
Nesta combinação de Categoria Administrativa e Grande Região, apenas nove dos cursos não receberam conceito ( SC ).
Os demais foram avaliados com conceito 1 ( 14 cursos ), conceito 2 ( 225 cursos ), conceito 4 ( 81 cursos ) e conceito 5 ( 40 cursos ).
Entre os 49 cursos em instituições públicas na região Sudeste, a categoria modal foi a que recebeu conceito 4 ( 16 cursos ).
Os demais foram avaliados com os conceitos 2 ( oito cursos ), 3 ( dez curso ) e 5 ( 15 cursos ).
Na região Sudeste, nesta categoria, nenhum curso ficou sem conceito ( SC ) ou recebeu conceito 1.
As instituições privadas concentraram 285 dos 336 cursos participantes da região Sul, 84,8 \% do total regional.
Desses, 148 ficaram com conceito 3, o conceito modal.
Nesta combinação de Categoria Administrativa e Grande Região, três cursos ficaram sem receber conceito ( SC ), enquanto nenhum recebeu conceito 1.
Os demais foram avaliados com conceito 2 ( 80 cursos ), conceito 4 ( 43 cursos ) e conceito 5 ( 11 curso ).
As instituições públicas na região Sul participaram com 51 cursos ( 15,2 \% ), dos quais nenhum ficou sem conceito.
O conceito modal também foi 3, com 20 cursos.
Além disso, os demais foram avaliados com conceito 1 ( um curso ), conceito 2 ( três cursos ), conceito 4 ( 19 cursos ) e conceito 5 ( oito cursos ).
Na região Centro-Oeste, 143 dos 171 cursos participantes eram de instituições privadas ( 83,6 \% em termos regionais ).
Destes, 73 concentraram -se no conceito 2, conceito modal.
Os demais receberam conceito 1 ( sete cursos ), conceito 3 ( 51 cursos ), conceito 4 ( oito cursos ) e conceito 5 ( quatro cursos ).
Nenhum curso ficou sem conceito ( SC ) nesta categoria.
Dos 28 cursos de instituições públicas, 12 foram avaliados no conceito 3, conceito modal.
Os demais receberam conceito 1 ( um curso ), conceito 2 ( seis cursos ), conceito 4 ( cinco cursos ) e conceito 5 ( quatro cursos ).
Nesta região, nenhum curso de IES pública ficou sem conceito ( SC ). 5.3 CONCEITOS POR ORGANIZAÇÃO ACADÊMICA E POR GRANDE REGIÃO Na Tabela 5.3 encontra -se a distribuição dos conceitos atribuídos aos cursos participantes do ENADE/2012 na Área de Administração, por Organização Acadêmica, segundo as Grandes Regiões brasileiras.
Dos 1.554 cursos participantes, 1.010 eram oferecidos em Faculdades, 402 em Universidades e os demais 142 em Centros Universitários.
Esta distribuição corresponde a, respectivamente, 65,0 \%, 25,9 \% e 9,1 \% dos cursos.
De acordo com os dados apresentados, de todos os 96 cursos avaliados com conceito 5, 50 eram vinculados a Universidades.
Somente um curso em Universidades ficou sem conceito ( SC ).
Esse tipo de Organização Acadêmica teve o conceito 3 como modal, com 173 cursos ( 43,0 \% ).
Os demais cursos de Universidades avaliados receberam os conceitos 1 ( três cursos ), 2 ( 81 cursos ), 4 ( 94 cursos ) e conceito 5 ( 50 cursos, como já mencionado ).
Entre os cursos em Centros Universitários, o conceito modal também foi 3, com 79 cursos ( 55,6 \% ).
Neste tipo de organização acadêmica somente um curso ficou sem conceito ( SC ).
Os outros cursos neste tipo de Organização Acadêmica receberam os conceitos 1 ( três cursos ), 2 ( 37 cursos ), 4 ( 17 cursos ) e conceito 5 ( cinco cursos ).
Nas Faculdades, 12 dos 1.010 cursos ficaram sem conceito ( SC ) e 420 ( 41,6 \% ) receberam o mesmo conceito modal dos demais tipos de Organização Acadêmica, conceito 3.
Dos demais cursos neste tipo de Organização Acadêmica, 27 receberam conceito 1, 406 conceito 2, 104 conceito 4 e 41 conceito 5.
Considerando -se separadamente as regiões brasileiras, verifica -se que, na região Norte, as Universidades concentraram 15 dos 88 cursos participantes.
Dos cursos em Universidades, nenhum recebeu conceitos 1 ou 5.
Seis cursos ficaram alocados igualmente aos conceitos 2 e 4 ( três em \textbf{cada} ).
Oito receberam o conceito modal, 3, enquanto apenas um curso não recebeu conceito ( SC ).
Os Centros Universitários da região Norte foram representados por sete cursos, destes, seis receberam o conceito modal, 3, e um recebeu conceito 4.
Nenhum curso ficou sem conceito ( SC ) ou recebeu os conceitos 1, 2 e 5.
As Faculdades participaram com 66 cursos na região Norte.
Nenhum ficou sem conceito ( SC ) ou recebeu o conceito 5.
Os demais receberam conceito 1 ( três cursos ), 2 ( 32 cursos, conceito modal ), conceito 3 ( 28 cursos ) e conceito 4 ( outros três cursos ).
Na região Nordeste, as Universidades participaram com 65 dos 262 cursos na Área de Administração da região.
Nenhum dos cursos em Universidades no Nordeste ficou sem conceito ( SC ) ou recebeu conceito 1.
O conceito modal foi 3, com 33 cursos.
Os demais receberam conceito 2 ( seis cursos ), 4 ( 19 cursos ) e conceito 5 ( sete cursos ).
Os Centros Universitários contaram com seis cursos participantes na região Nordeste, quatro deles com conceito 3 e os dois restantes com conceitos 2 e 5 ( um curso em \textbf{cada} ).
As Faculdades foram representadas por 191 cursos na região Nordeste, dos quais apenas um ficou sem conceito ( SC ).
O conceito modal foi 2, com 87 cursos.
Dos cursos restantes, sete receberam conceito 1, 73 conceito 3, 17 conceito 4 e seis receberam conceito 5.
Na região Sudeste, as Universidades concentraram 152 dos 697 cursos da região.
Entre os cursos em Universidades na região, o conceito modal foi 3 com 51 cursos, e nenhum curso ficou sem conceito ( SC ) ou recebeu conceito 1.
Os demais cursos receberam os conceitos 2 ( 44 cursos ), 4 ( 31 cursos ) e 5 ( 26 cursos ).
Os Centros Universitários participaram com 91 cursos na região Sudeste, dos quais 47 obtiveram conceito modal, 3, e nenhum ficou sem conceito ( SC ).
Os demais receberam os conceitos 1 ( três cursos ), 2 ( 27 cursos ), 4 ( dez cursos ) e 5 ( quatro cursos ).
As Faculdades foram representadas por 454 cursos na região Sudeste, que se distribuíram nos conceitos 1 ( 11 cursos ), 2 ( 162 cursos ), 3 ( 191 cursos, conceito modal ), 4 ( 56 cursos ) e 5 ( 25 cursos ).
Nove cursos ficaram sem conceito ( SC ).
Dos 336 cursos da região Sul, 129 eram de Universidades, para os quais o conceito modal foi 3, com 63 cursos.
Nesse tipo de organização, nenhum dos cursos ficou sem conceito ( SC ) ou recebeu conceito 1.
Os demais receberam os conceitos 2 ( 19 cursos ), 4 ( 36 cursos ) e 5 ( 11 cursos ).
Os Centros Universitários da região Sul tiveram 15 dos 26 cursos participantes no conceito modal, 3.
Dez cursos receberam conceitos 2 e 4 ( cinco \textbf{cada} ), enquanto um curso ficou sem conceito ( SC ).
Foram 181 os cursos vinculados a Faculdades na região Sul, e aos conceitos 2 e 4 ( três em \textbf{cada} ).
Oito receberam o conceito modal, 3, enquanto apenas um curso não recebeu conceito ( SC ).
Os Centros Universitários da região Norte foram representados por sete cursos, destes, seis receberam o conceito modal, 3, e um recebeu conceito 4.
Nenhum curso ficou sem conceito ( SC ) ou recebeu os conceitos 1, 2 e 5.
As Faculdades participaram com 66 cursos na região Norte.
Nenhum ficou sem conceito ( SC ) ou recebeu o conceito 5.
Os demais receberam conceito 1 ( três cursos ), 2 ( 32 cursos, conceito modal ), conceito 3 ( 28 cursos ) e conceito 4 ( outros três cursos ).
Na região Nordeste, as Universidades participaram com 65 dos 262 cursos na Área de Administração da região.
Nenhum dos cursos em Universidades no Nordeste ficou sem conceito ( SC ) ou recebeu conceito 1.
O conceito modal foi 3, com 33 cursos.
Os demais receberam conceito 2 ( seis cursos ), 4 ( 19 cursos ) e conceito 5 ( sete cursos ).
Os Centros Universitários contaram com seis cursos participantes na região Nordeste, quatro deles com conceito 3 e os dois restantes com conceitos 2 e 5 ( um curso em \textbf{cada} ).
As Faculdades foram representadas por 191 cursos na região Nordeste, dos quais apenas um ficou sem conceito ( SC ).
O conceito modal foi 2, com 87 cursos.
Dos cursos restantes, sete receberam conceito 1, 73 conceito 3, 17 conceito 4 e seis receberam conceito 5.
Na região Sudeste, as Universidades concentraram 152 dos 697 cursos da região.
Entre os cursos em Universidades na região, o conceito modal foi 3 com 51 cursos, e nenhum curso ficou sem conceito ( SC ) ou recebeu conceito 1.
Os demais cursos receberam os conceitos 2 ( 44 cursos ), 4 ( 31 cursos ) e 5 ( 26 cursos ).
Os Centros Universitários participaram com 91 cursos na região Sudeste, dos quais 47 obtiveram conceito modal, 3, e nenhum ficou sem conceito ( SC ).
Os demais receberam os conceitos 1 ( três cursos ), 2 ( 27 cursos ), 4 ( dez cursos ) e 5 ( quatro cursos ).
As Faculdades foram representadas por 454 cursos na região Sudeste, que se distribuíram nos conceitos 1 ( 11 cursos ), 2 ( 162 cursos ), 3 ( 191 cursos, conceito modal ), 4 ( 56 cursos ) e 5 ( 25 cursos ).
Nove cursos ficaram sem conceito ( SC ).
Dos 336 cursos da região Sul, 129 eram de Universidades, para os quais o conceito modal foi 3, com 63 cursos.
Nesse tipo de organização, nenhum dos cursos ficou sem conceito ( SC ) ou recebeu conceito 1.
Os demais receberam os conceitos 2 ( 19 cursos ), 4 ( 36 cursos ) e 5 ( 11 cursos ).
Os Centros Universitários da região Sul tiveram 15 dos 26 cursos participantes no conceito modal, 3.
Dez cursos receberam conceitos 2 e 4 ( cinco \textbf{cada} ), enquanto um curso ficou sem conceito ( SC ).
Foram 181 os cursos vinculados a Faculdades na região Sul, e dois destes ficaram sem conceito ( SC ).
De todos os cursos, 90 receberam conceito 3, o conceito modal.
Os demais receberam os conceitos 1 ( um curso ), 2 ( 59 cursos ), 4 ( 21 cursos ) e 5 ( oito cursos ).
Na região Centro-Oeste, 41 dos 171 cursos eram de Universidades.
Nesse tipo de organização, nenhum curso ficou sem conceito ( SC ) e o conceito modal foi 3, com 18 cursos.
Os outros cursos obtiveram os conceitos 1 ( três cursos ), 2 ( nove cursos ), 4 ( cinco cursos ) e 5 ( seis cursos ).
Os Centros Universitários da região Centro-Oeste contaram com 12 cursos, dos quais quatro receberam o conceito 2, sete o conceito 3, conceito modal, e um curso o conceito 4.
Nenhum curso ficou sem conceito ( SC ) ou recebeu os conceitos 1 ou 5.
Dos 118 cursos em Faculdades na região Centro-Oeste, nenhum ficou sem conceito ( SC ), cinco receberam conceito 1, 66 conceito 2, 38 conceito 3, sete conceito 4 e dois cursos conceito 5.

% matched lemmas: cada
\end{textsample}
